\documentclass{beamer}
\usefonttheme{professionalfonts}
\mode<presentation>{
  \usetheme{Luebeck}
  \usecolortheme{whale}
  \useoutertheme{infolines}
}
\setbeamertemplate{navigation symbols}{}
% [Beamer stop slide numbering after some slides - TeX - LaTeX Stack Exchange](https://tex.stackexchange.com/questions/609624/beamer-stop-slide-numbering-after-some-slides)
\setbeamertemplate{page number in head/foot}[appendixframenumber]
\usepackage[english]{babel}
\usepackage{csquotes}
\usepackage{mathtools}
\usepackage{unicode-math}
\usepackage{tikz}
\usepackage[absolute,overlay]{textpos}
\usetikzlibrary{backgrounds,fit,calc,patterns,positioning,arrows.meta,quotes,overlay-beamer-styles,cd}
\usepackage{forest}
\usepackage[conor]{../agda}
\usepackage{catchfilebetweentags}
\usepackage{amssymb}
\theoremstyle{remark}
\newtheorem{notation}[theorem]{Notation}
\newtheorem{remark}[theorem]{Remark}
\newtheorem{proposition}[theorem]{Proposition}

\usepackage[backend=biber,style=authoryear,maxcitenames=4,maxbibnames=10, url=false]{biblatex}

\addbibresource{../works.bib}

\usepackage{qrcode}
\usepackage{newunicodechar}
\DeclarePairedDelimiter{\expi}{\lceil}{\rceil}
\newunicodechar{⇒}{\ensuremath{\Rightarrow}}
%\DeclareUnicodeCharacter{21D2}{\ensuremath{\Rightarrow}}
\newunicodechar{ℰ}{\ensuremath{\mathscr{E}}}
\newunicodechar{₀}{\ensuremath{{}_{0}}}
\newunicodechar{ℕ}{\ensuremath{ℕ}}
\newunicodechar{∀}{\ensuremath{\forall}}
\newunicodechar{↦}{\ensuremath{↦}}
\newunicodechar{∘}{\ensuremath{∘}}
\newunicodechar{♯}{\ensuremath{_{\#}}}
\newunicodechar{ᵢ}{\ensuremath{_i}}
\newunicodechar{≔}{\ensuremath{≔}}
\newunicodechar{≈}{\ensuremath{\approx}}
\newunicodechar{₁}{\ensuremath{{}_{1}}}
\newunicodechar{₂}{\ensuremath{{}_{2}}}
\newunicodechar{λ}{\ensuremath{\lambda}}
\newunicodechar{≺}{\ensuremath{<}}
\newunicodechar{∈}{\ensuremath{\in}}
\newunicodechar{⊥}{\ensuremath{\bot}}
\newunicodechar{ε}{\ensuremath{\varepsilon}}
\newunicodechar{↓}{\ensuremath{\downarrow}}
%% Invisible characters for otherwise clashing definitions.
\newunicodechar{⇃}{\ensuremath{\!}}
\newunicodechar{⇂}{\ensuremath{\!}}
\newunicodechar{↿}{\ensuremath{\!}}
\newcommand{\sle}{\sqsubseteq} %sort \le
\newcommand{\sgt}{\sqsupset} %sort \ge
\newcommand{\inl}{\ensuremath{\mathsf{inl}}}
\newcommand{\inr}{\ensuremath{\mathsf{inr}}}
\newcommand{\id}{\ensuremath{\mathsf{id}}}
\newcommand{\sort}{\mathsf{sort}}
\newunicodechar{⊒}{\ensuremath{\sgt}}
\newunicodechar{⊑}{\ensuremath{\sle}}
\newunicodechar{ₗ}{\ensuremath{_l}}
\newunicodechar{ᵣ}{\ensuremath{_r}}
\newunicodechar{⁻}{\ensuremath{^-}}
\newunicodechar{¹}{\ensuremath{^1}}
\newunicodechar{∷}{\ensuremath{\mathbin{::}}}
\newunicodechar{┌}{|\ensuremath{\lceil}}
\newunicodechar{┐}{\ensuremath{\rceil}|}
\newunicodechar{₁}{\ensuremath{_1}}
\newunicodechar{₂}{\ensuremath{_2}}
\newunicodechar{→}{\ensuremath{\to}}
\newunicodechar{⸴}{\ensuremath{,}}
\newunicodechar{ℬ}{\ensuremath{\mathcal{B}\!}}
\newunicodechar{⊔}{\ensuremath{\sqcup}}
\newunicodechar{⊓}{\ensuremath{\sqcap}}
\newunicodechar{⊎}{\ensuremath{\uplus}}
\newunicodechar{≤}{\ensuremath{≤}}
\newunicodechar{▸}{\ensuremath{\blacktriangleright}}
\newfontfamily\cjkfont{NotoSansCJK-Regular.ttc}[FontIndex=0]
% Use \symbol, not the literal よ, inside the definition.
% newunicodechar makes よ active, so writing よ there again would recurse forever.
\newunicodechar{よ}{{\text{\cjkfont\symbol{"3088}}}}
%%%%
% elide "on obj/on morphs" distinction in code presentation
\newunicodechar{₀}{}


\tikzset{input/.style={circle,draw=red!50,fill=red!20,thick,
inner sep=0pt,minimum size=1cm}, output/.style={circle,draw=blue!50,fill=blue!20,thick,
inner sep=0pt,minimum size=1cm}}

% [beamer - Depth first step by step construction on forest tree - TeX - LaTeX Stack Exchange](https://tex.stackexchange.com/questions/489718/depth-first-step-by-step-construction-on-forest-tree?rq=1)
\tikzset{
    invisible/.style={opacity=0,text opacity=0},
    visible on/.style={alt=#1{}{invisible}},
    alt/.code args={<#1>#2#3}{%
      \alt<#1>{\pgfkeysalso{#2}}{\pgfkeysalso{#3}} % \pgfkeysalso doesn't change the path
    },
    %label visible on/.code={\tikzset{every label/.append style={visible on=#1}}}
}
\forestset{
  visible on/.style={
    for tree={
      /tikz/visible on={#1},
      edge+={/tikz/visible on={#1}}}}
  }


\forestset{
  ADT/.style={
    for tree={
      font={\fontsize{9.5pt}{9.5pt}\ttfamily},
      % https://www.dickimaw-books.com/latex/novices/html/fontstyle.html
    }
  },
}

\tikzset{
  algorithm/.style={
    algorithm line/.style={
      inner sep=0pt,
    },
    algorithm line label/.style={
    },
    algorithm split/.style={
      fill=red!30!white,
      draw=red!30!black,
    },
    striped line/.style={
      draw=none,
      pattern=north east lines,
      pattern color=##1,
    },
    boxed line/.style={
      draw=##1,
      line width=2pt,
      inner sep=-1pt,
      %pattern=crosshatch dots,
      %pattern color=##1,
    },
    line-orange/.style={
      draw=orange!70!black,
      pattern=north east lines,
      pattern color=orange!70!white,
    },
    line-blue/.style={
      draw=blue!70!black,
      pattern=north west lines,
      pattern color=blue,
    },
    line-yellow/.style={
      draw=yellow!70!black,
      pattern=north east lines,%crosshatch dots,
      pattern color=yellow!70!white,
    },
    line-green/.style={
      draw=green!70!black,
      pattern=north west lines,%horizontal lines,
      pattern color=green!70!white,
    },
    line-violet/.style={
      draw=purple!70!black,
      pattern=north east lines,
      pattern color=purple!70!white,
    },
  },
}


\usepackage{newunicodechar}
\usepackage{graphicx}

\title[Intrinsically Correct Total FP]{Intrinsically correct \& total functional programming using guarded recursion}
\author[Alexandru, Schaefer, New]{
  Cass Alexandru\inst{1} \and Steven Schaefer\inst{2} \and Max New\inst{2}
}
\institute[R(PT)U \& UM]{%
\inst{1}RPTU Kaiserslautern-Landau \& Radboud University Nijmegen
\and
\inst{2} University of Michigan
}
\date{2026-10-27}
\renewcommand<>{\emph}[1]{{\only#2{\bfseries}#1}}
\renewcommand{\hom}[3]{\ensuremath{#1\left(#2,#3\right)}}
\newcommand{\ihom}[2]{\ensuremath{\left[#1,#2\right]}}
\begin{document}
\begin{frame}
  \titlepage
\end{frame}

\begin{frame}
  \frametitle{Overview}
  \begin{itemize}
  \item Recent progress in intrinsically correct total functional programming \parencite{alexandruIntrinsicallyCorrectAlgorithms2026}
  \item Achieved with the novel notion of a \emph{well founded functor}
  \item \emph{Gap}: similarities to \emph{locally contractive functors} \parencite{birkedalFirstStepsSynthetic2012}
  \item \emph{Question}: Can we replace wff with lcf? \emph{Answer}: Yes
  \item Same application, stronger foundation
  \end{itemize}
\end{frame}

\begin{frame}
  \frametitle{Technical Overview}
  \begin{itemize}
  \item Locally contractive functors give rise to \emph{hylomorphisms}, which describe total functional algorithms
  \item To instantiate lcf, we need an \emph{enrichment} of our category (indexed datatypes for intrinsically correct programs, corresponding to \emph{families}) in \emph{presheaves}
  \item We go there and back again: Having instantiated the machinery, we expose a simpler interface for guarded recursion which uses only the \emph{family enrichment of families}
  \end{itemize}
\end{frame}

\section{Motivation}

\begin{frame}[fragile]
  \frametitle{QuickSort as a Coalgebra-to-Algebra Morphism}
\begin{columns}
  \begin{column}[t]{.7\textwidth}
    \begin{tikzcd}[ampersand replacement=\&,
every arrow/.append style={mapsto},column sep=2cm]
 {(i_{≤ p}, p, i_{> p})} \arrow[r, "\mathsf{sort} \times \mathrm{id}\times \mathsf{sort}"] \& {(\mathsf{sort}(i_{≤ p}), p, \mathsf{sort}(i_{> p}))} \arrow[d, "\mathsf{concat}"] \\
ip \arrow[u,"\mathsf{partition}"] \arrow[r,"\mathsf{sort}"] \&  {\mathsf{sort}(i_{≤ p})\, p\, \mathsf{sort}(i_{> p})}
\end{tikzcd}
\end{column}\hfill
\begin{column}[t]{.3\textwidth}
\begin{tikzcd}
F_0I \arrow[r, "F_1h"]\ar[phantom, start anchor=center, end anchor=center]{dr}[description]{\refstepcounter{equation}(\theequation)\label{eq:square}} & F_0O \arrow[d, "a"] \\
I \arrow[u,"c"] \arrow[r,"h"] &  O
\end{tikzcd}
  \end{column}
\end{columns}
\begin{itemize}
\item for \(\mathsf{partition}\): \(F ≔ B \) for \(B X ≔ 1 + X × A × X\), \(\quad I ≔ A^*\).
\item A c2a-morphism \(h\) is called a \emph{hylomorphism} if it is the \emph{unique} c2a-morphism making \eqref{eq:square} commute.
\item NB: In a non-total programming language, \texttt{h := c ; F1 h ; a} \emph{is} already a definition. We want to recover a similar definition scheme in a \emph{total} language.
\end{itemize}
\end{frame}

\begin{frame}
  \frametitle{Pars pro Toto : Total from Partial Correctness}
  \begin{itemize}
  \item The type \(A^* → 1 + A^* × A × A^*\) for \(\mathsf{partition}\) is \emph{too coarse!}
  \item As \cite{AlexandruCRW25}: Move to \(\mathsf{Set}/{ℳ\!A}\)
  \end{itemize}
  \vspace{\baselineskip}%
\[
  \overline{B} \begin{pmatrix}
    X \\
    f
  \end{pmatrix}
  ≔
  \begin{pmatrix}
    1 \\
    * ↦ \emptyset
  \end{pmatrix}
  +
  \begin{pmatrix}
    \{(l,p,r) ∈ X × A × X \mid \dots\} \\ %
    (l, p, r) ↦ f(l) \uplus \{p\} \uplus f(r)
  \end{pmatrix}
\]
\begin{itemize}
\item Refined type of partition: \(\mathsf{Set}/{ℳ\!A}[(A^*, \mathsf{toBag}),\overline{B}(A^*, \mathsf{toBag})]\)
\item Note: The multiset indices in functorial positions are smaller than the input index: \(f(l),f(r) < f(l) \uplus \{p\} \uplus f(r)\)
\end{itemize}
\end{frame}

\begin{frame}[fragile]
  \frametitle{Approach}
  \begin{itemize}
  \item Observation: \enquote{indices in functorial positions are smaller than the input index}
  \end{itemize}
  \begin{equation*}
\begin{tikzcd}
F_0I \arrow[r, "F_1h"]\ar[phantom, start anchor=center, end anchor=center]{dr}[description]{(1)} & F_0O \arrow[d, "a"] \\
I \arrow[u,"c"] \arrow[r,"h"] &  O
\end{tikzcd}    
  \end{equation*}
  \begin{itemize}
  \item \(⇒\) we could construct a solution to (1) \enquote{as an inhabitant of the homset using well founded recursion}!
\item Next: How this is \emph{made formal}. We begin with moving from Hom-\emph{sets} to Hom-\emph{objs} for some category with the requisite structure for constructing an inhabitant using well founded recursion
  \end{itemize}
\end{frame}

% we can encode the correctness by indexing over the contents of the list…
\section{Background}

\begin{frame}[fragile]
  \frametitle{Enriched Categories \& Enrichment}
 Fix a set of objects \(C_0\) and a monoidal category \((V,⊗,𝟙)\).
\begin{definition}[Enriched category]
  \begin{itemize}
  \item for every \(x, y \colon C_0\) the internal Hom \(\ihom{X}{Y}\);
  \item an internal identity morphism \(\mathrm{id}_X \colon \hom{V}{𝟙}{\ihom{X}{X}}\) \& sequencing operation \(\hom{V}{\ihom{X}{Y}⊗\ihom{Y}{Z}}{\ihom{X}{Z}}\);
  \item internal counterparts to unitality \& associativity coherences
  \end{itemize}
\end{definition}
\begin{itemize}
\item For an enriched category \(C'\) there is the \emph{underlying} plain (i.e. \(\mathsf{Set}\)-enriched) category with \(ΓC'[X,Y] ≔ \hom{V}{𝟙}{\ihom{X}{Y}}\)
\item But we're starting \emph{out} with a plain category!
\end{itemize}
\begin{definition}[Enrichment of a category \(C\)]
  \begin{itemize}
  \item same data \& laws as above;
  \item \(C ≅ ΓC'\), in particular \(\expi{-}\colon \hom{C}{X}{Y} ≅ \hom{V}{𝟙}{\ihom{X}{Y}}\)
  \end{itemize}
\end{definition}
\end{frame}

\begin{frame}
  \frametitle{Families enriched in Families}
  \begin{itemize}
  \item Families \((\mathsf{Set}^{C_0})\) are Cartesian monodial, inheriting that structure pointwise from \(\mathsf{Set}\)
  \item \(𝟙 = \prod_{c ∈ C_0} 𝟙_{\mathsf{Set}}\)
  \item \(\ihom{X}{Y} c = \hom{\mathsf{Set}}{X c }{Y c}\)
  \item \(\hom{\mathsf{Set}^{C_0}}{X }{Y} = \prod_{c ∈ C_0}\hom{\mathsf{Set}}{X c }{Y c} ≅ \prod_{c ∈ C_0}\hom{\mathsf{Set}}{𝟙_{\mathsf{Set}}}{\hom{\mathsf{Set}}{X c }{Y c}} = \hom{\mathsf{Set}^{C_0}}{𝟙 }{\ihom{X}{Y}} \)
  \end{itemize}
\end{frame}

\begin{frame}
  \frametitle{Enrichment of a Functor}
\begin{definition}[Enrichment of a Functor]
  Given:
  \begin{itemize}
  \item categories \(C, D\) with enrichments \(EC = (\ihom{-}{-},\expi{-})\) and \(ED = (\ihom{-}{-}',\expi{-}')\) in monoidal category \(V\);
  \item a functor \(F \colon C → D\)
  \end{itemize}
   An \emph{enrichment} of \(F\) in \(EC,ED\) is an internal functorial action
   \[F_1' \colon \hom{V}{\ihom{X}{Y}}{\ihom{F_0 X}{F_0 Y}'}\]
   \vspace{-1\baselineskip{}}
   \begin{itemize}
   \item obeying internal counterparts to preservation of id and comp.;
   \item and such that \(F_1\) and \(F_1'\) \emph{agree}: (\(\expi{-}\colon \hom{C}{X}{Y} ≅ \hom{V}{𝟙}{\ihom{X}{Y}}\))
   \end{itemize}
  \[∀ X, Y : C_0,\, f : \hom{C}{X}{Y}.\ \expi{F_1 f}' = \expi{f} ; F_1' \]
\end{definition}  
\end{frame}

\begin{frame}
  \frametitle{Locally Contractive Functors}
  \begin{definition}
    Given:
  \begin{itemize}
  \item \(V\) a \enquote{model of guarded recursive terms}
  \item Categories \(C,D\) with enrichments \((EC,\expi{-}), (ED,\expi{-}')\) in \(V\)
  \item Functor \(F \colon C → D\)
  \item An enriched functor \(F' \colon ▸EC → ED\) with \(F_0\) as object map
  \end{itemize}
    \(F\) is \emph{locally contractive} if it has as enrichment the sequence of enrichments \(EC \xrightarrow{\mathsf{next}} ▸EC \xrightarrow{F'} ED\), i.e.
\[∀ x, y : C_0,\, f : \hom{C}{x}{y}.\ \expi{F_1 f}' = \expi{f} ; \mathsf{next} ; F_1' \]
\end{definition}
\begin{itemize}
\item NB: This is actually a corollary of the normal definition, for \(F\) with no preexisting enrichment.
\end{itemize}
\end{frame}

\begin{frame}
  \frametitle{the boundary of the representable: \(\mathsf{PSh}(A,<)\) as a model of guarded recursive terms}
  Need : \(\mathsf{▸}, \mathsf{next}, \mathrm{Löb}\)
  \begin{itemize}
  \item \(↓ i = \{ j ≤ i\}. (↓ i) ≔ \hom{A}{–}{i} = よ i\). \(↡ i ≔ \{ j < i \}\)
  \item \(▸ X i = \mathrm{lim}_{j<i} X j = \hom{\mathsf{PSh}(A)}{∂よ i }{X}\)
  \item \(\mathsf{next}\colon \hom{\mathsf{Endo}(\mathsf{PSh}(A))}{\mathsf{Id}}{▸}\)
    %, \(\mathsf{next} (X , \mathsf{restrX}_{{j ≤ i} \colon \hom{\mathsf{Set}}{X i}{X j}}) i \colon \hom{\mathsf{Set}}{(↡ i) j }{X j }\)
  \item \(\mathsf{next}\) given by functoriality of the presheaf
  \item \(\mathrm{Löb}:\) for \(H \colon \hom{\mathsf{PSh}(A)}{▸ \ihom{X}{Y}}{\ihom{X}{Y}}\) \(\mathrm{Löb}(H) \colon \hom{\mathsf{PSh}(A)}{𝟙}{\ihom{X}{Y}}\) s.t. \(\mathrm{Löb}H = \mathrm{Löb}H ; \mathsf{next} ; H\)
  \item \(\mathrm{Löb}\) given by well founded induction
  \end{itemize}
\end{frame}


\begin{frame}[fragile]
  \frametitle{Locally Contractive Functors give Hylomorphisms}
  \framesubtitle{Proof sketch} % Throw out 90% of the symbols
  \begin{columns}
\begin{column}[t]{.3\textwidth}
\begin{tikzcd}
F_0I \arrow[r,dashed,"1?"]\ar[phantom, start anchor=center, end anchor=center]{dr}[description]{(1)} & F_0O \arrow[d, "a"] \\
I \arrow[u,"c"] \arrow[r,dashed,"0?"] &  O
\end{tikzcd}
\end{column}
\begin{column}{.7\textwidth}
  Constraints:
  \begin{enumerate}
  \item \(1? = F_1 0?\)
  \item \(0? = c ; 1? ; a\)
  \end{enumerate}
\end{column}
\end{columns}
\begin{itemize}
\item Note that in \eqref{eq:square}, the functional \(H(h) ≔ c ; h ; a\) has type \((F_0 I → F_0 O) → (I → O)\).
\item Local contractivity of \(F\) gives us \(F_1' \colon ▸ (I → O) → (F_0I → F_0O) \).
\item Constraint 1 is given by \(F_1'\) agreeing with \(F_1\).
\item \(\mathrm{Löb} \colon (▸ \ihom{I}{O} → \ihom{I}{O}) → \ihom{I}{O}\) gives us \(0?\) by \(\mathrm{Löb} (F_1' ; H)\).
\item Constraint 2 follows from using \(\mathrm{Löb}\).
\end{itemize}
\end{frame}

\section{Instantiation}

\begin{frame}
  \frametitle{Change of Base of Enrichment}
Let \((M,*,I)\) and \((N,⊗,J)\) be monoidal categories.
\begin{definition}[Lax monoidal functor]
  A functor \(F \colon M → N\) is \emph{lax monoidal} if it has comparison arrows:
  \begin{itemize}
\item \(ε \colon \hom{N}{J}{F I}\), and
\item \(μ_{X,Y} \colon \hom{N}{F x ⊗ F Y}{F (x * y)}\) agreeing with the unitors and associators of \(M\) and \(N\).
\end{itemize}
\end{definition}

\begin{proposition}[Change of Base]
  Let \(EC= (\ihom{-}{-},\dots)\) be an enriched category for \(C_0\) in \(M\). Let \(F\colon M → N\) be a lax monoidal functor. We can define an enriched category \(F^*(EC)\) for \(C_0\) in \(N\) called the base-change of \(EC\) along \(F\) with \(\ihom{X}{Y} = F (\ihom{X}{Y})\).
\end{proposition}
\end{frame}

\begin{frame}
  \frametitle{Families Enriched in Presheaves}
  \begin{itemize}
  \item We already saw families enriched in families
  \item To enrich them in Presheaves, we need a lax monoidal functor \(\mathsf{Fam}(A) → \mathsf{PSh}(A)\)
  \item The functor that does the job is the \emph{cofree presheaf functor}: \(U ⊣ \mathsf{Cofree}\), \(\mathsf{Cofree}(B)(x) ≔ \prod_y \hom{\mathsf{Set}}{\hom{A}{y}{x}}{B(y)} = \prod_{y ≤ x} B(y)\)
  \item lax monoidal structure derives from RAPL
  \end{itemize}
\end{frame}

\begin{frame}
  \frametitle{Families are Enough: There and Back Again}
  % we have later at home
  \begin{itemize}
  \item We can define a version of \(▸\) for families by \(▸' ≔ U ∘ ▸ ∘ \mathsf{Cofree}\)
  \item Have the enrichment of families in families with internal homs \(ℰ\)
  \item for lcf, the contractive action of \(F\) lives in \(\hom{\mathsf{Psh}(A)}{▸ (\mathsf{Cofree}\, \ihom{X}{Y})}{\mathsf{Cofree}\, \ihom{FX}{FY}}\)
  \item \(≅\langle  \) transpose this homset along \(U ⊣ \mathsf{Cofree}\)! \(\rangle\)
  \item \(\hom{\mathsf{Fam}(A)}{U ▸ \mathsf{Cofree}\, \ihom{X}{Y}}{\ihom{FX}{FY}} = \hom{\mathsf{Fam}(A)}{▸' \ihom{X}{Y}}{\ihom{FX}{FY}}\)
  \item The agreement condition can be similarly ported to the setting of families only
  \item This yields an expression of lcf for enrichment in families
  \item Transposed along \(U ⊣ \mathsf{Cofree}\), this is equivalent to the usual lcf
  \end{itemize}
\end{frame}

\section{Outlook}

% \begin{frame}
%   \frametitle{Related Work}
%   \begin{itemize}
%   \item Advantages of lcf over wff:
%     \begin{itemize}
%     \item Consolidates the research landscape
%     \item Some technical improvements that make the category theory agree exactly with the provided library for writing intrinsically correct \& total algorithms while requiring fewer theorems as premise
%     \item By instantiating an existing notion we:
%       \begin{itemize}
%       \item make this new application more accessible to those already familiar with lcf;
%       \item profit from using instantiations of existing results.
%       \end{itemize}
%     \end{itemize}
%   \end{itemize}
% \end{frame}
    %   \begin{itemize}
    % \item Novel instantiation of lcf on \(𝒞^A\), for \((A, <)\) a well founded order
    % \item all coalgebras for a lcf are recursive (corollary of known result) \(⇒\) usable for same applications as wff
    % \item instantiation achieved by an \emph{enrichment} of \(𝒞^A\) in \(\mathsf{PSh}(A,<)\)
    % \item Families are enough: Enrichment in families and a new condition \enquote{lc-Fam} sufficient for lcf
    % \end{itemize}

\begin{frame}
  \frametitle{Conclusion}
  \begin{itemize}
  \item Novel application of locally contractive functors to total functional programming
  \item Can be used \enquote{behind the scenes} as a theoretical underpinning while the programmer need only work in the setting of indexed datatypes (families)
  \item Illustrates an interesting novel usecase for lcf in a category which is not self-enriched
  \end{itemize}
\end{frame}

\begin{frame}
  \frametitle{TODOs}
  \begin{itemize}
  \item Actually see which existing results about lcf are \emph{actually} useful to instantiate for applications
  \item In particular: Can we use composition theorems to reduce client obligations? Famously local contractivity as is closed under composition with enriched functors
  \item See whether solving domain equations using lcf is also useful in this setting
  \item Technical: Currently we're switching back between enriched categories \& enrichments because some lax monoidal functors do not preserve underlying categories. Need to make this more elegant
  \end{itemize}
\end{frame}

\end{document}
